\begin{table}[H]
\centering
\small
\caption{Análisis tornado: impacto de cada variable en el VAN}
\label{tab:tabla9_tornado}
\begin{tabularx}{\linewidth}{>{\raggedright\arraybackslash}X>{\raggedright\arraybackslash}X>{\raggedleft\arraybackslash}X>{\raggedleft\arraybackslash}X>{\raggedleft\arraybackslash}X}
\toprule
\textbf{Variable} & \textbf{Rango evaluado} & \textbf{VAN desfavorable (S/)} & \textbf{VAN favorable (S/)} & \textbf{Amplitud (S/)} \\
\midrule
Tipo de cambio & -20 \% / +20 \% & -195,218 & 6,149,577 & 6,344,795 \\
Precio FOB & -20 \% / +20 \% & -195,218 & 6,149,577 & 6,344,795 \\
Rendimiento & -20 \% / +20 \% & 976,359 & 4,996,839 & 4,020,480 \\
COK & 16 \% / 8 \% & 1,900,751 & 4,437,132 & 2,536,380 \\
Costo variable & +20 \% / -20 \% & 1,840,373 & 4,146,720 & 2,306,347 \\
Fracción exportable & 76 \% / 95 \% & 1,428,713 & 3,550,331 & 2,121,619 \\
Costo fijo & +20 \% / -20 \% & 2,035,819 & 3,946,813 & 1,910,993 \\
Inversión & +20 \% / -20 \% & 2,542,752 & 3,445,213 & 902,461 \\
\bottomrule
\end{tabularx}
\par\vspace{2pt}\parbox{\linewidth}{\footnotesize\textit{Nota.} VAN base = S/ 2,993,982. Elaboración propia.}
\end{table}
